\section{Protocolo de consentimiento informado}
\label{sec:consentimiento-kano}
\begingroup
\small
\setlength{\parskip}{3pt}

\noindent\textbf{Identificación institucional y académica.} El presente protocolo corresponde al equipo MRS, integrado por Mendoza Palma Allan Jeremy, Rochina Rochina Freddy David y Sánchez Centeno Roselyn Andreina, estudiantes de Ingeniería de Software de la Universidad Técnica Estatal de Quevedo. Se desarrolla en el marco de la asignatura Ingeniería de Requerimientos (ISR-401), bajo la docencia de Gleiston Cicerón Guerrero Ulloa, PhD, en la actividad FORM-SIM-01 del período SPA 2026--2027.

\noindent\textbf{Propósito.} Se solicita la participación en la aplicación de un cuestionario basado en el modelo de Kano para priorizar ocho características del dominio ``Unidad educativa: matrícula y calificaciones''. La actividad tiene fines académicos y formativos, no genera calificación para quien participa y sirve de preparación para la actividad individual TA-IND-02.

\noindent\textbf{Qué se solicita.} Cada persona responde, para cada una de las ocho características, dos preguntas: una sobre cómo se sentiría si la característica estuviera presente y otra sobre cómo se sentiría si estuviera ausente. Cada respuesta se elige entre cinco opciones. La duración estimada es de diez a quince minutos por persona; esta estimación queda sujeta a confirmación al definirse la modalidad de aplicación.

\noindent\textbf{Voluntariedad y retiro.} La participación es voluntaria. La persona puede omitir cualquier pregunta, interrumpir la sesión o retirarse sin consecuencias académicas ni personales. Si retira su consentimiento, el equipo deja de utilizar las respuestas correspondientes y gestiona su exclusión del análisis.

\noindent\textbf{Seudonimizado y tratamiento de la información.} En los archivos del repositorio público las personas participantes se identifican solo con los códigos P01 a P05. No se publican nombres, cédulas, teléfonos, correos electrónicos, audios, videos ni imágenes identificables. El uso de códigos constituye una seudonimización: reduce la exposición directa, pero no garantiza anonimato absoluto si existen datos adicionales que permitan establecer una vinculación. La información se emplea exclusivamente para el análisis académico de la actividad y se trata conforme a los principios aplicables de la normativa ecuatoriana de protección de datos personales~\cite{lopdp2021} y del Código de Ética y Conducta Profesional de la ACM~\cite{acm2018}. Este modelo no interpreta disposiciones legales concretas.

\noindent\textbf{Almacenamiento y acceso.} Los formularios firmados y cualquier evidencia que contenga datos personales se conservan únicamente en una carpeta institucional privada vinculada a cuentas institucionales de la UTEQ, con acceso restringido al equipo y al docente. En el repositorio público permanece solo el modelo sin completar, el archivo de respuestas seudonimizado y el manifiesto de evidencias. La conservación se limita al tiempo necesario para la evaluación académica. El presente modelo no implica aprobación institucional previa ni garantía jurídica absoluta.

\noindent\textbf{Manifestación del consentimiento.} Antes de iniciar, quien aplica el instrumento explica las condiciones anteriores y solicita la aceptación de forma verbal o escrita. La aceptación permite comenzar el cuestionario; la negativa impide su aplicación. Si se utiliza el formulario escrito, se completa en la copia privada y no se publica.

\noindent\textbf{Registro de la aplicación.} Para la aplicación del 27 de noviembre de 2026, Rochina comunicó que las personas participantes aceptaron el consentimiento. El repositorio dispone únicamente de las respuestas seudonimizadas exportadas: cinco formularios con los códigos P01, P04, P04, P02 y P05. No se cuenta con copias de los originales ni con registros que permitan confirmar si la aceptación fue verbal o escrita, la duración efectiva de cada sesión o los roles por código. Por tal motivo, el manifiesto público no incorpora filas ni huellas SHA-256 y no se publica información personal de quienes participaron.

\vspace{4pt}
\begin{center}
\renewcommand{\arraystretch}{1.8}
\begin{tabular}{@{}>{\bfseries\raggedright\arraybackslash}p{0.31\linewidth}>{\raggedright\arraybackslash}p{0.60\linewidth}@{}}
Nombre de quien participa & \rule[-2pt]{0.95\linewidth}{0.4pt} \\
Fecha & \rule[-2pt]{0.95\linewidth}{0.4pt} \\
Modalidad del consentimiento & Verbal\quad \fbox{\rule{0pt}{1.6ex}\rule{1.6ex}{0pt}}\quad Escrita\quad \fbox{\rule{0pt}{1.6ex}\rule{1.6ex}{0pt}} \\
Firma (solo si es escrita) & \rule[-2pt]{0.95\linewidth}{0.4pt} \\
Quien aplica el instrumento & \rule[-2pt]{0.95\linewidth}{0.4pt} \\
\end{tabular}
\end{center}
\endgroup
